JHQ screens routed candidates with primary codes before reading residuals
for a smaller survivor set \cite{jhq}. For a query $q$, let $U(q)$ be
the candidates in its selected IVF lists. The refinement multiplier sets
\begin{equation}
 c(\alpha)=\max\{k,\lceil\alpha k\rceil\},\qquad
 C_\alpha(q)=\operatorname{Top}^{\rm primary}_{\min\{c(\alpha),|U(q)|\}}U(q).
 \label{eq:6-budget}
\end{equation}
Here $C_\alpha$ denotes valid survivors. The implementation reserves
$c(\alpha)$ slots and can pad an underfull pool with sentinels. The
analysis assumes at least $k$ valid candidates. Raising $\alpha$ expands
refinement capacity within the visited lists; raising $n_{\rm probe}$
changes the visited universe. Neither operation recovers information
lost by the other automatically.

A large fixed $\alpha$ may repeatedly refine candidates that do not
change the final top-$k$. A small value can exclude a useful neighbour
before residual scoring. Which tradeoff is worthwhile depends on the
query population, primary reconstruction, list coverage and refined-score
ordering. One global default therefore need not suit every dataset--probe
configuration. The evaluated grid spans $c=20$ to $2{,}000$ at $k=10$.
Changing it reuses the same trained index and stored residuals.
Lowering the multiplier reduces residual rows accessed, rather than
releasing the resident residual representation. Survivor buffers grow
when capacity is insufficient and retain that capacity afterward, so a
smaller budget is not a proportional reduction in allocated memory.

\begin{figure}[tbp]
\centering
% Logical work boundary, not a measured latency or pruning diagram.
\begin{tikzpicture}[
 font=\footnotesize\linespread{1}\selectfont,
 box/.style={draw=ThesisNavy!65,rounded corners=2pt,align=center,
 text width=2.15cm,minimum height=1.5cm,inner sep=3pt},
 arr/.style={-{Latex[length=2mm]},draw=ThesisNavy!80,thick}]
\node[box] (route) at (0,0) {Probe IVF lists\\Fixed $n_{\rm probe}$};
\node[box] (primary) at (3,0) {Score $U(q)$\\Primary codes};
\node[box,fill=ThesisNavy!7] (select) at (6,0)
 {Primary top-$c$\\Keep $C_\alpha(q)$};
\node[box,fill=ThesisNavy!7] (refine) at (9,0)
 {Refine $c(\alpha)$ slots\\Residual rows};
\node[box,fill=ThesisNavy!7] (final) at (12,0)
 {Final top-$k$\\Return $k$ IDs};
\draw[arr] (route) -- (primary);
\draw[arr] (primary) -- (select);
\draw[arr] (select) -- (refine);
\draw[arr] (refine) -- (final);
\node[align=center,text width=5.1cm] at (1.5,1.4)
 {Fixed routed universe and number of primary scores};
\node[align=center,text width=7.6cm] at (9,1.4)
 {Budget-dependent selection, residual scoring and final selection};
\draw[densely dashed,ThesisNavy!45] (4.5,-0.9) -- (4.5,1.1);
\node[align=center,text width=13.8cm] at (6,-1.45)
 {$c(\alpha)=\max\{k,\lceil\alpha k\rceil\}$ slots;
 valid survivors are capped by $|U(q)|$.\\
 The stored model and residual codes are reused across budgets.};
\end{tikzpicture}

\caption[Work affected by the refinement budget]{The refinement budget changes work after primary-score
production. With a fixed query, index and probe count, it leaves the
routed candidate universe unchanged. The highlighted stages depend on
survivor capacity; this is a logical work decomposition, not a measured
latency breakdown. The scan kernel also performs budget-dependent
admission and compaction. Padding applies when fewer than $c$ candidates exist.}
\label{fig:6-budget-scope}
\end{figure}

An explanatory latency decomposition is
\begin{equation}
 T(\alpha)=T_{\rm base}+T_{\rm select}(c(\alpha))+
 T_{\rm refine}(c(\alpha))+T_{\rm final}(c(\alpha)).
 \label{eq:6-cost-model}
\end{equation}
For a fixed batch and route, $T_{\rm base}$ includes transfer, rotation,
routing and primary-score production; selection, residual scoring and
final selection depend on survivor capacity (\cref{fig:6-budget-scope}).
Buffer capacities, occupancy and resource thresholds can change in steps,
so this work model does not predict a monotone runtime. The measured
small-budget plateau is evaluated in \cref{sec:7-budget-trace}; the policy
minimises accepted budget rather than a noisy timing rank.


Each calibration trial still runs rotation, probing and primary scoring
before refinement. Reusing the index avoids rebuilding or retraining;
it does not cache every intermediate query result across budget trials.
The cost of calibration must therefore count full searches, not merely
the difference between two residual-scoring workloads.

The policy selects one budget for an index, routing configuration, $k$
and query workload. It uses neither a trained per-query predictor nor
an online choice for each arriving query. Scan-group scheduling remains
the separate mechanism of \cref{sec:5-5}; its execution path must still
be considered when interpreting budget calibration.

\paragraph{Position among adaptive ANN methods.}
The multiplier has the same candidate-pool role as FAISS
\texttt{IndexRefine::k\_factor} and cuVS's refinement ratio
\cite{faiss_refine,cuvs_refinement}. FAISS auto-tuning explores
speed--accuracy operating points using a ground-truth criterion
\cite{faiss_autotune}. Li et al. learn per-query termination conditions
from search features \cite{li_termination}; DARTH predicts recall during
execution to stop at a requested target \cite{darth}.
Here the decision calibrates one unlabelled workload budget for reuse,
using reference-output coverage rather than a trained recall predictor.
The final scorer remains a compressed reconstruction, whereas the cuVS
refinement example reranks source vectors. These distinctions separate
the work parameter, its selection mechanism and the scored representation.
